\subsection{I-38. Two centered pairs with isolated invertible affine factors}
\label{I-38}

\paragraph{Statement.}
The conclusion and sample bound of I-37 also hold for a width-at-most-three
affine univariate matrix word with the following larger promise:
after constant factors have been absorbed, no two consecutive
variable-dependent factors both have generic rank three. Thus
arbitrarily many generic-rank-three affine factors may occur,
provided a singular factor separates each two of them. The order
remains unknown, the nodes are any distinct working-field elements,
and the conclusion holds in every characteristic.

Equivalently, for any nonzero slopes $c_i$ and distinct nodes $a_i$,
the rational map $R_{2,i}=tc_i+z_1/(a_i-y_1)+z_2/(a_i-y_2)$
hits the stated class. The diagonal rescaling (I.99) then proves
ordinary common-center $t\mathbf1+\mathrm{SV}_2$ nonvanishing.
We prove the rational formulation.

\paragraph{1. The one-pole center dominates every three coordinates.}
Let $\alpha_i=tc_i+z/(a_i-y)$. Pairwise dominance is (I.92).
For distinct $i,j,k$, solve for $t,z$ from $X_i,X_j$ over $E(y)$.
The third coordinate has the form
\[
 X_k=A(y)X_i+B(y)X_j,
 \qquad A,B\in E(y).
 \tag{I.101}
\]
The two coefficients cannot both be constant. If they were,
comparison of the $z$ coefficients would give
$1/(a_k-y)=A/(a_i-y)+B/(a_j-y)$, which the distinct poles
exclude. Hence the right side of (I.101) is a nonconstant
rational function of $y$ over $E(X_i,X_j)$, and is transcendental
over that field. Thus the three coordinates are algebraically
independent. This proves dominance in every characteristic,
without a Jacobian argument. It also proves the claim over every
fixed extension of the coefficient field.

\paragraph{2. The representation promise survives irreducible reduction.}
Assume a nonzero annihilator and apply I-37's prime-kernel and
disjoint-variable factor argument. Specializing the variables
outside an irreducible factor makes the intervening matrix
factors constant. Absorb each constant block into the next
remaining factor, and absorb the final block into the terminal
column. All remaining factors are affine and their ranks cannot
increase.

The no-adjacent-rank-three promise also persists. If two adjacent
remaining factors both had rank three, both corresponding
original factors would have had rank three. The original promise
puts a singular original factor between them. It has now been
specialized into the intervening constant block, so that block
has rank at most two. Absorbing it into the later factor makes
that factor singular, a contradiction. This reasoning does not
require any specialized singular factor to keep its original
rank. A zero specialization cannot occur in the realized nonzero
scalar word.

We may therefore assume that the input is irreducible and depends
on every remaining variable. As in I-37, a strict internal
rank-one interval has a polynomial outer factorization and
splits the scalar into two nonconstant factors. Every strict
internal interval consequently has generic rank two or three.
If there are at most three remaining variables, part~1 already
hits the scalar. Assume $n\ge4$ below.

\paragraph{3. One pole preserves every internal interval rank.}
Each generic-rank-three factor is invertible at its univariate
center, because its nonzero determinant is univariate and the
center is transcendental. Its second compound is likewise
invertible there.

Consider a generic-rank-two interval. In its second compound
word, a singular original factor becomes a rank-one factor,
and every intervening rank-three factor becomes invertible.
There is at most one such intervening factor between consecutive
rank-one factors. Write a nonzero rank-one factor as $p_iq_i/b_i$
using a nonzero univariate pivot. The connecting scalar between
two successive rank-one factors is therefore either
\[
 q_i(X_i)p_j(X_j)
 \quad\hbox{or}\quad
 q_i(X_i)C_2(A_k(X_k))p_j(X_j).
 \tag{I.102}
\]
It depends on at most three distinct coordinates. Since the
original compound interval product is nonzero, every connecting
scalar is nonzero. Part~1 preserves all these scalars and all
univariate pivots. The remaining rank-one endpoint vectors are
nonzero at their centers, and any full-rank prefix or suffix is
invertible there. The compound interval product stays nonzero,
so the original interval still has rank two. A generic-rank-three
interval consists entirely of invertible factors and retains
rank three. Thus all strict internal interval ranks are preserved
simultaneously. The proof requires preservation of their ranks;
it does not assert preservation of every individual two-minor.

Absorb $A_1,A_n$ into the affine endpoint row and column. There
is an internal rank-two factor: otherwise the at-least-two
internal positions would be consecutive rank-three factors.
The first internal rank-two factor is at position two or three.
The endpoint row times the prefix through that factor is a
nonzero polynomial row in at most three variables, and remains
nonzero at the center by part~1. All subsequent constant
intervals have rank at least two, so propagation of its active
two-dimensional row space is injective. Hence the left endpoint
has a nonzero active projection for the full constant interior
word. The reversed argument at the last internal rank-two
factor gives a nonzero active projection of the right endpoint.

\paragraph{4. Compatible $I/P$ constants and endpoint shears.}
Over the field $K=\overline E(t,z,y)$, the constant interior word
has full-product rank two, and every subinterval rank two or
three. We give the compatible-space argument explicitly.
Use row-vector orientation. Choose a two-dimensional initial
space $U_0$ complementary to the left kernel of the full
constant product, and propagate it through the constants to
spaces $U_i$. Each propagation is an isomorphism on $U_i$.
Choose a complementary line $V_m$ at the final cut. For the
suffix product $T_i$ from cut $i$ to the end, put
\[
 V_i=\{v:vT_i\in V_m\}.
\]
If $T_i$ has rank two, its row image equals $U_m$, so
$V_i=\ker T_i$ is a line. If it has rank three, the inverse
image of $V_m$ is a line. In either case $U_i\cap V_i=0$,
and each constant factor sends the corresponding passive line
into the next one. Active bases transported through the word
make every active block $I_2$. The passive block is zero at
rank-two factors and a nonzero scalar at rank-three factors;
rescale the successive passive bases to make the latter scalar
one. Thus all interior factors have the form
\[
 A_i(\alpha_i+Y_i)=C_i+Y_iD_i,
 \qquad C_i\in\{P,I_3\},\quad P=\operatorname{diag}(I_2,0),
 \tag{I.103}
\]
with $C_i=P$ precisely at the rank-two positions. No consecutive
interior constants are both $I_3$.

The nonzero active endpoint vectors pair to
$f(\boldsymbol\alpha)=0$. Normalize them to $e_1^{\mathsf T}$
and $e_2$ by one common active change of basis. The endpoint
passive components must also be removed. Use the upper shear
$L$ from I-37 at every cut of the initial string of $I_3$
constants, up to the first $P$. Then
$L^{-1}I_3L=I_3$ along that string and $L^{-1}P=P$ at its end,
while the initial passive row component is killed. Similarly
use the lower shear $R$ from I-37 at every cut of the terminal
string of $I_3$ constants after the last $P$:
$PR=P$ and $R^{-1}I_3R=I_3$. This kills the passive column
component. If the first and last $P$ are the same factor,
$L^{-1}PR=P$ permits both shears simultaneously. Thus all
constants in (I.103) remain as stated, and
\[
 f(\boldsymbol\alpha+\mathbf Y)
 =(e_1^{\mathsf T}+Y_1p)
   \prod_{i=2}^{n-1}(C_i+Y_iD_i)
   (e_2+Y_nq).
 \tag{I.104}
\]

The pole extraction of I-37 gives $f_i(\boldsymbol\alpha)=0$
for every $i$. Therefore $p_2=q_1=(D_i)_{12}=0$. At every
singular factor, the coefficient of $Y_i$ in
$\det(P+Y_iD_i)=0$ also gives $(D_i)_{33}=0$. An invertible
factor may have a nonzero passive derivative entry; it is not
silently discarded.

\paragraph{5. Two tangent directions kill the bandwidth-two Hessian.}
Set $h_i=1/(a_i-y)$. Differentiate the identities
$f_i(\boldsymbol\alpha)=0$ with respect to $z$ and $y$.
Equivalently, extract their first coefficients under independent
parameter increments. Since $z\ne0$ in the rational function
field, this gives in every characteristic
\[
 H(\boldsymbol\alpha)\mathbf h=0,
 \qquad H(\boldsymbol\alpha)\mathbf h^{\circ2}=0,
 \tag{I.105}
\]
where $\mathbf h^{\circ2}=(h_1^2,\ldots,h_n^2)$. The diagonal
of $H$ is zero by multilinearity.

If $i<j$ and a singular interior position lies strictly between
them, the intervening constants contain $P$. Their product
is $P$, which removes the passive coordinate, while the active
derivative blocks are lower triangular. Hence $H_{ij}=0$.
If there is no such position, all intervening constants are
$I_3$, and $H_{ij}$ is exactly the product of the state-one-to-
passive entry at position $i$ and the passive-to-state-two
entry at position $j$, with the endpoint versions $p_3$ and
$q_3$ when appropriate. There are no alternative active
contributions because all $(1,2)$ entries vanish.

Since no two consecutive internal positions are invertible,
every gap of length at least three contains a singular
intermediate position. Thus $H_{ij}=0$ whenever $|i-j|>2$.
In the first row, at most two future entries remain. For two
indices $j\ne k$, the corresponding coefficient determinant
in (I.105) is
\[
 \det\begin{pmatrix}h_j&h_k\\h_j^2&h_k^2\end{pmatrix}
 =h_jh_k(h_k-h_j)\ne0.
 \tag{I.106}
\]
Both entries are therefore zero. If only one entry remains,
$h_j\ne0$ suffices. By symmetry the earlier entries in the
next row are now zero; apply the same argument to its at most
two future entries. Induction gives $H=0$.

Finally expand (I.104) into state paths, as in I-37. A path
cannot move directly from active state one to state two.
Before its first arrival at state two, take its last departure
from state one to the passive state, or its initial passive
endpoint entry if it has not yet visited state one. Every
intermediate passive stay must occur at an invertible factor:
a singular factor has both passive constant and passive
derivative entries zero. The product of the departure and
arrival entries is the corresponding Hessian entry $H_{ij}$,
and is zero. Intermediate passive stays only multiply this
zero product. The same argument covers an output through the
passive terminal component; an output through active state one
is excluded by $q_1=0$. Therefore every output path is zero,
contradicting the assumed nonzero scalar. This proves the
rational two-pair assertion and hence the ordinary centered
two-pair assertion.\hfill$\square$

\paragraph{Sampling, provenance and the remaining gate.}
The exact field and interpolation requirements are unchanged
from I-37: for total degree at most $D\le n$, the sample count
is (I.100), over a working field with at least
$\max\{n,D+1,(n-1)D+1\}$ distinct elements, or an extension
with the explicitly stated evaluation convention. No input
order or representation is reconstructed by the sampler.

The extra proof ingredients are the function-field proof of
triple dominance, the preserved adjacency promise under
irreducible specialization, the rank-one compound connectors
through one invertible factor, and the two tangent kernel
identities. Purity independently suggested the use of
$\mathbf h,\mathbf h^{\circ2}$ instead of slope-dependent
minors, avoiding an invalid Jacobian shortcut in characteristic
two. No historical priority claim is made.

Two or more consecutive invertible factors remain outside
this proof. Their compound connectors involve at least four
coordinates, so three-coordinate dominance alone does not
preserve their rank; even at a rank-preserved center, the
Hessian can have more than two future entries in a row. The
larger semiseparable Hessian structure and the consequences
of the assumed two-pair identity for possible one-pole rank
loss remain the specific next gate.
